% =====================================================================
%  THÈME 4 — Réactif limitant et avancement
%  Niveau FACILE — Exercices
% =====================================================================
\documentclass[11pt,a4paper]{article}
\usepackage[utf8]{inputenc}
\usepackage[T1]{fontenc}
\usepackage{lmodern}
\usepackage[french]{babel}
\usepackage{amsmath,amssymb,mathtools}
\usepackage{icomma}
\usepackage[version=4]{mhchem}
\usepackage{siunitx}
\sisetup{output-decimal-marker={,},per-mode=symbol}
\usepackage{xcolor}
\usepackage{tikz}
\usepackage[most]{tcolorbox}
\usepackage{enumitem}
\usepackage{array}
\usepackage{fancyhdr}
\usepackage[a4paper,margin=2.1cm,headheight=26pt,headsep=8pt]{geometry}

\definecolor{primary}{RGB}{142,93,168}
\definecolor{primarydark}{RGB}{82,45,108}
\definecolor{excol}{RGB}{0,135,90}
\newcommand{\nq}[1]{n_{\text{#1}}}

\pagestyle{fancy}\fancyhf{}
\fancyhead[L]{\small\textcolor{primarydark}{\textbf{Fiche 5} — Thème 4 : Limitant}}
\fancyhead[R]{\small\textcolor{excol}{\textsc{Facile}}}
\fancyfoot[C]{\small\thepage}
\renewcommand{\headrulewidth}{0.8pt}
\renewcommand{\headrule}{\hbox to\headwidth{\color{primary}\leaders\hrule height \headrulewidth\hfill}}

\newtcolorbox[auto counter]{exo}[1][]{enhanced,breakable,boxrule=1pt,arc=3pt,
  left=3mm,right=3mm,top=2.4mm,bottom=2.4mm,colback=excol!5,colframe=excol,
  fonttitle=\bfseries,coltitle=white,
  title={Exercice~\thetcbcounter\ \ifx\relax#1\relax\relax\else\textmd{--- \itshape #1}\fi}}

\setlength{\parindent}{0pt}
\setlength{\parskip}{3pt}

\begin{document}

\begin{center}
{\LARGE\bfseries\color{primarydark}Thème 4 — Réactif limitant \& avancement}\\[1pt]
{\large\color{excol}\textbf{Niveau Facile}}\\[2pt]
{\itshape Objectif : identifier le réactif limitant par le critère du plus petit rapport $n/\nu$.}\\
{\color{primary}\rule{0.62\linewidth}{1pt}}
\end{center}

\begin{exo}[la synthèse générique]
Pour $\ce{2 A + 3 B -> C + 6 D}$, on mélange $\nq{A,0}=\SI{0,100}{\mole}$ et $\nq{B,0}=\SI{0,200}{\mole}$.
\textbf{a)} Calcule le rapport $n/\nu$ pour \ce{A} et pour \ce{B}.\quad
\textbf{b)} Quel est le réactif limitant ? Quelle est la valeur de $\xi_{\max}$ ?
\end{exo}

\begin{exo}[synthèse de l'ammoniac]
Pour $\ce{N2 + 3 H2 -> 2 NH3}$, on introduit \SI{2,0}{\mole} de \ce{N2} et \SI{9,0}{\mole} de \ce{H2}.
Quel est le réactif limitant ?
\end{exo}

\begin{exo}[formation de l'eau]
Pour $\ce{2 H2 + O2 -> 2 H2O}$, on dispose de \SI{5,0}{\mole} de \ce{H2} et \SI{2,0}{\mole} de \ce{O2}.
Quel est le réactif limitant ?
\end{exo}

\begin{exo}[le piège : « plus petite quantité » $\neq$ limitant]
On fait réagir le zinc avec l'acide chlorhydrique : $\ce{Zn + 2 HCl -> ZnCl2 + H2}$.
On introduit \SI{0,50}{\mole} de \ce{Zn} et \SI{0,80}{\mole} de \ce{HCl}.
\textbf{a)} Quel réactif est présent en plus grande quantité de matière ?\quad
\textbf{b)} En appliquant le critère $n/\nu$, quel est le réactif \emph{limitant} ? Commente.
\end{exo}

\begin{exo}[vrai ou faux ? notions de limitant et d'excès]
Pour une réaction où un réactif est limitant, indique si chaque affirmation est \textbf{vraie}
ou \textbf{fausse}, et justifie brièvement :
\begin{enumerate}[label=\arabic*.,leftmargin=1.8em,topsep=1pt,itemsep=1pt]
  \item « Le réactif limitant est toujours celui dont la quantité de matière est la plus faible. »
  \item « À la fin de la réaction, il reste du réactif en excès. »
  \item « Le réactif limitant est entièrement consommé. »
  \item « Tous les réactifs sont entièrement consommés. »
\end{enumerate}
\end{exo}

\end{document}
